\documentclass[10pt,a4paper]{amsart}
\usepackage[margin=19mm]{geometry}
\usepackage{amsmath,amssymb,mathtools}
\usepackage[scr=rsfs,cal=euler]{mathalfa}
\usepackage{xfrac}
\usepackage[unicode,hidelinks,bookmarks=false]{hyperref}
\newcommand{\ie}{i.e.\ }
\newcommand{\cf}{cf.\ }
\newcommand{\resp}{respectively\ }
\newcommand{\Subsection}{\subsection}
\providecommand{\nobreakdash}{\nobreak\hbox{-}}
\setlength{\emergencystretch}{2em}
\multlinegap0pt
\usepackage{amssymb}
\usepackage{mathtools}
\usepackage{bm}
\usepackage{algorithm}\usepackage{algpseudocode}
\usepackage{float}
\usepackage{xcolor}
\newtheorem{lemma}{Lemma}
\renewcommand{\Im}{\operatorname{Im}}
\newcommand{\ii}{i}
\title{Adaptive local representations for Helmholtz Trefftz discontinuous Galerkin methods}
\author{Shelvean Kapita}
\date{Source: arXiv:2609.12282v1 (CC BY 4.0)}
\begin{document}
\maketitle

\noindent\emph{Source and licence.} Adaptive local representations for Helmholtz Trefftz discontinuous Galerkin methods. Version arXiv:2609.12282v1 (CC BY 4.0), distributed by arXiv under the Creative Commons Attribution 4.0 International licence, http://creativecommons.org/licenses/by/4.0/, as recorded in the arXiv licence metadata for this exact version and shown on the abstract page https://arxiv.org/abs/2609.12282v1. Full licence notice: \texttt{licenses/arxiv-2609.12282-CC-BY-4.0.txt}.\par\smallskip

\noindent\textit{Editorial scope.} Counted unit: the article's lemma lem:pw-direction-lipschitz (source0.tex lines 608--641). The article's complete original proof of it is delivered with it (source0.tex lines 642--656, one \texttt{proof} environment of 15 lines). No mathematics is written, repaired or re-derived: the statement and the proof are the article's own lines, in the article's own notation, and no retained line is substituted or re-flowed.  No further source-local context is needed: every cross-reference of every retained line resolves inside the two delivered spans.  Nothing was omitted from any retained span, so the package carries no omission record.  No retained line contains a citation, so the wrapper carries no bibliography and no bibliography entry.  No retained line contains a figure environment, an image-inclusion command or drawing code, so no image is delivered and the package declares no asset dependency.  Editorial formatting only, in addition: an AMS \texttt{amsart} preamble that restates the article's own declarations and macros used by the retained lines, namely the environments \texttt{lemma} on separate counters, together with the standard packages those lines need.  The retained environments print in this document as Lemma 1 in order of appearance, whereas the article prints the same items with the numbering of its own full document.  The complete native source of the article as distributed in the arXiv e-print for this exact version is bundled unchanged as \texttt{sources/210095/source0.tex}.
\begin{lemma}[Complex-angle Lipschitz bounds]\label{lem:pw-direction-lipschitz}
Let $K\subset B(x_K,h_K)$ and $|\Im z|,|\Im w|\le\eta_{\max}$.  Define
\[
 D_\eta=\sqrt{\cosh(2\eta_{\max})},
 \qquad E_K=e^{\kappa h_K\sinh\eta_{\max}},
\]
and introduce the augmented element-boundary norm
\begin{equation}\label{eq:boundary-plus-norm}
 \|v\|_{\partial K,+}^2
 :=\int_{\partial K}\big(\kappa|v|^2+\kappa^{-1}|\nabla v|^2\big)\,ds.
\end{equation}
Then
\begin{align}
 \|\psi_z^K-\psi_w^K\|_{H^1(K)}
 &\le C_{H^1,K}|z-w|,\label{eq:pw-H1-lipschitz}\\
 \|\psi_z^K-\psi_w^K\|_{\mathrm{tr},K}
 &\le C_{\mathrm{tr},K}|z-w|,\label{eq:pw-tr-lipschitz}\\
 \|\psi_z^K-\psi_w^K\|_{\partial K,+}
 &\le C_{\partial,K}|z-w|,\label{eq:pw-boundary-lipschitz}
\end{align}
where one may take
\begin{align}
 C_{H^1,K}
 &=|K|^{1/2}\kappa D_\eta E_K
 \Big(h_K^2+(1+\kappa h_KD_\eta)^2\Big)^{1/2},\label{eq:CH1K}\\
 C_{\mathrm{tr},K}
 &=\Big(2\pi h_K\kappa D_\eta^2E_K^2
 [ (\kappa h_K)^2+(1+\kappa h_KD_\eta)^2]\Big)^{1/2},\label{eq:CtrK}\\
 C_{\partial,K}
 &=\Big(|\partial K|\,\kappa D_\eta^2E_K^2
 [ (\kappa h_K)^2+(1+\kappa h_KD_\eta)^2]\Big)^{1/2}.\label{eq:CboundaryK}
\end{align}
For $\eta_{\max}=0$ these reduce to the corresponding real-angle bounds.  The factor $E_K=e^{\kappa h_K\sinh\eta_{\max}}$ is intrinsic: strongly evanescent directions are exponentially amplified when continued across a patch of radius comparable with $h_K$.  This is the reason for imposing a bounded complex-angle strip rather than allowing arbitrarily large $|\Im z|$.
\end{lemma}

\begin{proof}
Along the segment joining $z$ and $w$,
\[
 \partial_z\psi_z=\ii\kappa d'(z)\cdot(x-x_K)\psi_z,
 \qquad
 \partial_z\nabla\psi_z=\ii\kappa d'(z)\psi_z+\ii\kappa d(z)\partial_z\psi_z.
\]
On the strip, $\|d(z)\|_2,\|d'(z)\|_2\le D_\eta$ and $|\psi_z|\le E_K$.  Hence, uniformly for $x\in B(x_K,h_K)$,
\[
 |\partial_z\psi_z|\le \kappa D_\eta h_KE_K,
 \qquad
 |\partial_z\nabla\psi_z|\le \kappa D_\eta E_K(1+\kappa h_KD_\eta).
\]
Integration along the segment in the complex-angle plane gives the pointwise difference bounds.  Integration over $K$, over $\partial B(x_K,h_K)$, and over $\partial K$, respectively, yields \eqref{eq:pw-H1-lipschitz}--\eqref{eq:pw-boundary-lipschitz}.  On the circle, $|\partial_n v|\le|\nabla v|$, which gives \eqref{eq:CtrK}.
\end{proof}

\end{document}
